% =============================================================================
% Преамбула учебника «Цифровая схемотехника: от логических элементов до ПЛИС»
% Компиляция: xelatex main.tex (дважды, для оглавления и ссылок)
% =============================================================================

% --- Языки и шрифты (XeLaTeX) ---
\usepackage{fontspec}
\usepackage{polyglossia}
\setdefaultlanguage{russian}
\setotherlanguage{english}

% На Windows используются Times New Roman / Arial / Consolas,
% на Linux -- метрически совместимые Liberation и DejaVu Sans Mono.
\IfFontExistsTF{Times New Roman}
  {\def\MainFontName{Times New Roman}}{\def\MainFontName{Liberation Serif}}
\IfFontExistsTF{Arial}
  {\def\SansFontName{Arial}}{\def\SansFontName{Liberation Sans}}
\IfFontExistsTF{Consolas}
  {\def\MonoFontName{Consolas}}{\def\MonoFontName{DejaVu Sans Mono}}
\setmainfont{\MainFontName}[Ligatures=TeX]
\setsansfont{\SansFontName}[Ligatures=TeX]
\setmonofont{\MonoFontName}[Scale=MatchLowercase]
% polyglossia требует явно указать шрифты для кириллицы
\newfontfamily\cyrillicfont{\MainFontName}[Ligatures=TeX]
\newfontfamily\cyrillicfontsf{\SansFontName}[Ligatures=TeX]
\newfontfamily\cyrillicfonttt{\MonoFontName}[Scale=MatchLowercase]

\usepackage{microtype}
\emergencystretch=3em

% --- Абзацы ---
\usepackage{indentfirst}
\usepackage{setspace}
\setstretch{1.15}
\setlength{\parindent}{1.25cm}
\setlength{\parskip}{3pt plus 1pt}

% --- Математика ---
\usepackage{amsmath}
\usepackage{amssymb}

% --- Поля ---
\usepackage{geometry}
\geometry{a4paper, top=22mm, bottom=22mm, left=25mm, right=20mm,
          headheight=15mm, footskip=12mm}

% --- Цвета ---
\usepackage[table]{xcolor}
\definecolor{FpgaBlue}{RGB}{14, 107, 78}
\definecolor{FpgaDark}{RGB}{14, 33, 27}
\definecolor{FpgaOrange}{RGB}{214, 158, 20}
\definecolor{FpgaGreen}{RGB}{84, 138, 38}
\definecolor{FpgaRed}{RGB}{200, 30, 45}
\definecolor{FpgaPurple}{RGB}{120, 40, 160}
\definecolor{FpgaLight}{RGB}{237, 244, 240}
\definecolor{FpgaGray}{RGB}{110, 125, 118}
\definecolor{BoardGreen}{RGB}{20, 90, 60}
\definecolor{ChipBlack}{RGB}{35, 35, 40}
\definecolor{CodeBack}{RGB}{247, 249, 248}
\definecolor{CodeComment}{RGB}{0, 128, 0}
\definecolor{CodeKeyword}{RGB}{0, 60, 170}
\definecolor{CodeString}{RGB}{170, 60, 0}

% --- Графика ---
\usepackage{etoolbox}
\usepackage{xstring}
\usepackage{graphicx}
\usepackage{float}
\usepackage{tikz}
\usetikzlibrary{arrows.meta, calc, positioning, shapes.geometric, shapes.misc,
  shapes.gates.logic.US, shapes.gates.logic.IEC, shapes.symbols,
  fit, backgrounds, decorations.pathreplacing, patterns, shadows, matrix}
\usepackage{tikz-timing}
\usepackage{pgfplots}
\pgfplotsset{compat=1.18}
\usepgfplotslibrary{groupplots}
\usepackage[siunitx, RPvoltages]{circuitikz}

% --- Заголовки, колонтитулы ---
\usepackage[automark, headsepline]{scrlayer-scrpage}
\clearpairofpagestyles
\ihead{\headmark}
\ohead{\textbf{\textcolor{FpgaBlue}{ЦИФРОВАЯ СХЕМОТЕХНИКА}}}
\cfoot{\pagemark}
\setkomafont{pageheadfoot}{\sffamily\small\color{FpgaDark}}
\setkomafont{pagenumber}{\sffamily\bfseries\color{FpgaBlue}}
\setkomafont{chapter}{\sffamily\LARGE\bfseries\color{FpgaDark}}
\setkomafont{section}{\sffamily\Large\bfseries\color{FpgaDark}}
\setkomafont{subsection}{\sffamily\large\bfseries\color{FpgaBlue}}
\setkomafont{subsubsection}{\sffamily\normalsize\bfseries\color{FpgaOrange}}
\setkomafont{paragraph}{\sffamily\bfseries\color{FpgaDark}}
\setkomafont{part}{\sffamily\Huge\bfseries\color{FpgaBlue}}
\renewcommand*{\chapterformat}{%
  \colorbox{FpgaBlue}{\color{white}\,\thechapter\,}\enskip}
\RedeclareSectionCommand[beforeskip=2.2ex plus .4ex, afterskip=1.1ex plus .2ex]{section}
\RedeclareSectionCommand[beforeskip=1.8ex plus .3ex, afterskip=.8ex plus .2ex]{subsection}

% --- Списки и таблицы ---
\usepackage{enumitem}
\setlist{itemsep=1pt, parsep=0pt, topsep=3pt}
\setlist[itemize]{label=\textcolor{FpgaBlue}{\textbullet}, leftmargin=*}
\setlist[enumerate]{label=\textcolor{FpgaBlue}{\bfseries\arabic*.}, leftmargin=*}
\setlist[description]{font=\bfseries\color{FpgaDark}}
\usepackage{booktabs}
\usepackage{tabularx}
\usepackage{multirow}
\usepackage{longtable}
\usepackage{caption}
\captionsetup{font=small, labelfont={bf,color=FpgaBlue}}
\usepackage{subcaption}

% --- Листинги кода (minted v3: нужен Pygments и ключ -shell-escape; папку сборки задаёт TEXMF_OUTPUT_DIRECTORY) ---
\usepackage{minted}
\usemintedstyle{friendly}
\setminted{fontsize=\small, breaklines=true, tabsize=4, baselinestretch=1.0}

% --- Цветные блоки ---
\usepackage[most]{tcolorbox}
\tcbuselibrary{minted}
% =============================================================================
% Копирование кода из PDF без искажений.
%
% Проблема: при копировании листинга из PDF теряются отступы, дефис
% подменяется типографским, а длинные строки рвутся. Решение (перенесено из
% основного учебника personal/lessons/tex/styles/preamble.tex): перед
% вёрсткой читаем исходный файл и прикрепляем к каждой строке её точный
% текст через /ActualText (пакет accsupp). Внешний вид не меняется.
% Нумерация строк в листингах не используется, чтобы номера не попадали в
% копируемый текст.
% =============================================================================
\usepackage{accsupp}
\makeatletter
\def\ACCSUPP@bdc{\special{pdf:code \ACCSUPP@span BDC}}
\def\ACCSUPP@emc{\special{pdf:code EMC}}
\makeatother

\ExplSyntaxOn
\seq_new:N  \g__codeacc_hex_seq
\int_new:N  \l__codeacc_index_int
\int_new:N  \l__codeacc_pending_int
\ior_new:N  \g__codeacc_ior
\bool_new:N \g__codeacc_active_bool
\tl_new:N   \l__codeacc_line_tl
\str_new:N  \l__codeacc_hex_str

% Store one raw source line as UTF-16BE hex (or an empty entry for a line
% that will not produce any glyph on the page).
\cs_new_protected:Npn \__codeacc_store_line:n #1
  {
    \tl_set:Nx \l__codeacc_line_tl { \tl_trim_spaces:n {#1} }
    \tl_if_empty:NTF \l__codeacc_line_tl
      {
        \int_incr:N \l__codeacc_pending_int
        \seq_gput_right:Nn \g__codeacc_hex_seq { }
      }
      {
        \str_set_convert:Nnnn \l__codeacc_hex_str {#1} { } { utf16be/hex }
        \seq_gput_right:Nx \g__codeacc_hex_seq { \str_use:N \l__codeacc_hex_str }
        \int_zero:N \l__codeacc_pending_int
      }
  }

% Read the whole source file once, before \inputminted typesets it.
\cs_new_protected:Npn \__codeacc_load:n #1
  {
    \seq_gclear:N \g__codeacc_hex_seq
    \bool_gset_false:N \g__codeacc_active_bool
    \file_if_exist:nT {#1}
      {
        \int_zero:N \l__codeacc_pending_int
        \ior_open:Nn \g__codeacc_ior {#1}
        \ior_str_map_inline:Nn \g__codeacc_ior { \__codeacc_store_line:n {##1} }
        \ior_close:N \g__codeacc_ior
        \bool_gset_true:N \g__codeacc_active_bool
      }
  }

% fvextra calls \FancyVerbFormatLine more than once per line (once inside a
% measuring \sbox), so the source line is looked up through fancyvrb's own
% FancyVerbLine counter instead of a counter of our own.
\cs_new_protected:Npn \__codeacc_wrap_line:n #1
  {
    \bool_if:NTF \g__codeacc_active_bool
      {
        \int_set:Nn \l__codeacc_index_int { \value{FancyVerbLine} }
        \str_set:Nx \l__codeacc_hex_str
          { \seq_item:Nn \g__codeacc_hex_seq { \l__codeacc_index_int } }
        \str_if_empty:NTF \l__codeacc_hex_str
          { #1 }
          {
            \exp_args:Nx \BeginAccSupp
              { method=hex,unicode,ActualText= \str_use:N \l__codeacc_hex_str }
            #1
            \EndAccSupp { }
          }
      }
      { #1 }
  }

% Оборачиваем \inputminted: так исправление работает и для \vfile, и для
% фрагментов vcode (tcolorbox выводит их через \inputminted из временного
% файла).
\cs_set_eq:NN \__codeacc_orig_inputminted:w \inputminted
\RenewDocumentCommand \inputminted { O{} m m }
  {
    \__codeacc_load:n {#3}
    \group_begin:
      \cs_set_eq:NN \FancyVerbFormatLine \__codeacc_wrap_line:n
      \__codeacc_orig_inputminted:w [#1] {#2} {#3}
    \group_end:
  }
\ExplSyntaxOff

\tcbset{codestyle/.style={enhanced, breakable, listing only, listing engine=minted,
  colback=CodeBack, colframe=FpgaBlue, boxrule=0pt, leftrule=2.2pt, arc=0pt,
  left=6pt, right=4pt, top=2pt, bottom=2pt, before skip=6pt, after skip=8pt,
  fonttitle=\ttfamily\small\bfseries, coltitle=FpgaDark, colbacktitle=FpgaLight,
  attach boxed title to top left={yshift=-1pt}, boxed title style={boxrule=0pt, arc=0pt}}}
% Фрагмент кода: \begin{vcode} ... \end{vcode}, язык по умолчанию Verilog
\newtcblisting{vcode}[1][verilog]{codestyle, minted language=#1,
  minted options={fontsize=\small, breaklines=true, tabsize=4}}
% Файл целиком: \vfile[язык]{путь}{подпись} (без номеров строк -- см. codecopy.tex)
\newtcbinputlisting{\vfile}[3][verilog]{codestyle, minted language=#1,
  minted options={fontsize=\small, breaklines=true, tabsize=4},
  listing file={#2}, title={#3}}
\newtcolorbox{infobox}[1][Справка]{
  enhanced, breakable, colback=FpgaLight, colframe=FpgaBlue, frame hidden,
  borderline west={3pt}{0pt}{FpgaBlue}, sharp corners=east, arc=1.5mm, boxrule=0pt,
  colbacktitle=FpgaBlue!14, coltitle=FpgaDark, fonttitle=\sffamily\bfseries, title={#1}}
\newtcolorbox{warnbox}[1][Внимание!]{
  enhanced, breakable, colback=orange!6, colframe=FpgaOrange, frame hidden,
  borderline west={3pt}{0pt}{FpgaOrange}, sharp corners=east, arc=1.5mm, boxrule=0pt,
  colbacktitle=FpgaOrange!14, coltitle=FpgaDark, fonttitle=\sffamily\bfseries, title={#1}}
\newtcolorbox{taskbox}[1][Задание]{
  enhanced, breakable, colback=green!4, colframe=FpgaGreen, frame hidden,
  borderline west={3pt}{0pt}{FpgaGreen}, sharp corners=east, arc=1.5mm, boxrule=0pt,
  colbacktitle=FpgaGreen!14, coltitle=FpgaDark, fonttitle=\sffamily\bfseries, title={#1}}
\newtcolorbox{ideabox}[1][Главная мысль]{
  enhanced, breakable, colback=purple!4, colframe=FpgaPurple, frame hidden,
  borderline west={3pt}{0pt}{FpgaPurple}, sharp corners=east, arc=1.5mm, boxrule=0pt,
  colbacktitle=FpgaPurple!14, coltitle=FpgaDark, fonttitle=\sffamily\bfseries, title={#1}}

% --- Прочее ---
\usepackage{siunitx}
\sisetup{output-decimal-marker={,}, group-separator={\,}, locale=FR}
\usepackage{hyperref}
\hypersetup{colorlinks=true, linkcolor=FpgaBlue, urlcolor=FpgaBlue,
  citecolor=FpgaBlue, pdftitle={Цифровая схемотехника: от логических элементов до ПЛИС},
  pdfauthor={Корягин С.А.}}

% Удобные команды
\newcommand{\code}[1]{\texttt{\small #1}}
\newcommand{\term}[1]{\textbf{\textcolor{FpgaDark}{#1}}}
\newcommand{\codepath}{../code}

% Общие стили TikZ
\tikzset{
  block/.style={draw=FpgaBlue, fill=FpgaLight, thick, rounded corners=2pt,
                minimum height=9mm, minimum width=22mm, align=center,
                font=\sffamily\small},
  hblock/.style={block, fill=FpgaBlue, text=white},
  oblock/.style={block, draw=FpgaOrange, fill=orange!10},
  gblock/.style={block, draw=FpgaGreen, fill=green!8},
  arr/.style={-{Stealth[length=2.5mm]}, thick, color=FpgaDark},
  wire/.style={thick, color=FpgaDark},
  bus/.style={line width=2pt, color=FpgaBlue},
  lbl/.style={font=\sffamily\scriptsize, color=FpgaDark},
}
\tikzset{
  gbase/.style={draw=FpgaDark, thick, fill=FpgaLight, logic gate inputs=nn,
                minimum width=12mm, minimum height=9mm, inner sep=0pt},
  gAND/.style={and gate US, gbase},
  gOR/.style={or gate US, gbase},
  gXOR/.style={xor gate US, gbase},
  gNAND/.style={nand gate US, gbase},
  gNOR/.style={nor gate US, gbase},
  gXNOR/.style={xnor gate US, gbase},
  ibase/.style={draw=FpgaDark, thick, fill=FpgaLight, logic gate inputs=nn,
                minimum width=8mm, minimum height=10mm, font=\sffamily\small},
  iAND/.style={and gate IEC, ibase},
  iOR/.style={or gate IEC, ibase},
  iXOR/.style={xor gate IEC, ibase},
  iNAND/.style={nand gate IEC, ibase},
  iNOR/.style={nor gate IEC, ibase},
  iXNOR/.style={xnor gate IEC, ibase},
  iNOT/.style={not gate IEC, draw=FpgaDark, thick, fill=FpgaLight,
               minimum width=8mm, minimum height=7mm, font=\sffamily\small},
  gNOT/.style={not gate US, draw=FpgaDark, thick, fill=FpgaLight,
               minimum width=9mm, minimum height=6mm, inner sep=0pt},
}

% --- Снимки экрана Yosis ---
\graphicspath{{../img/}}
\usepackage{adjustbox}
% \shot[ширина]{файл}{подпись}{метка}
\newcommand{\shot}[4][\textwidth]{%
  \begin{figure}[H]
  \centering
  \adjustbox{max width=#1, cfbox=FpgaGray!60 0.4pt 0pt}{\includegraphics[width=#1]{#2}}
  \caption{#3}\label{#4}
  \end{figure}}
